\begin{frame}{소수 극단 샘플이 업데이트를 지배한다}
  \begin{columns}[T]
    \begin{column}{0.60\textwidth}
      \centering
      \includegraphics[width=\linewidth]{ch11_1_ratio_variance.png}
      {\footnotesize 좌: $\mu$ 이동에 따른 $\operatorname{Var}[r]
      =e^{\mu^2}-1$ (로그 축). 우: $\mu=2$ 일 때 10,000개 샘플 중
      상위 1\%가 전체 가중치 합의 약 31\% 차지.}
    \end{column}
    \begin{column}{0.38\textwidth}
      {\small
      $\mu=2$ 로 10,000개 샘플:
      \begin{itemize}
        \item 가중치 평균은 1에 가까움 (0.91).
        \item \textbf{최댓값은 약 143}.
        \item \textbf{상위 1\%(100개)} 샘플이 전체 가중치 합에서
          약 \textbf{31\%} 차지.
      \end{itemize}
      }
      \vskip0.5em
      {\small 업데이트 방향을 정하는 것은 ``대부분의 샘플''이
      아니라 \textbf{소수의 극단 샘플} — 극단 샘플이 노이즈이므로
      업데이트가 노이즈에 지배당해 정책이 흔들린다.
      바로 11.2절 클리핑이 ``극단 샘플의 기여를
      $1+\epsilon$ 으로 잘라버리는'' 이유.}
    \end{column}
  \end{columns}
\end{frame}
